\documentclass{article} % Clase base sin tamaño

% Página horizontal
\usepackage[a4paper,landscape,margin=2cm]{geometry}
\usepackage{float} 
% Arial y tamaño exacto
\usepackage{fontspec}
\setmainfont{Arial}
\usepackage{xcolor}
% Para tamaño grande explícito
\usepackage{setspace}
\usepackage{graphicx}

\usepackage{amsmath}
\usepackage[thmmarks, amsmath, thref]{ntheorem}
\usepackage{amsfonts}
\usepackage{amssymb}
\usepackage{polyglossia}       % Paquete multilingüe moderno
\setmainlanguage{spanish}   

\newcommand{\bigtable}[1]{{\fontsize{30pt}{36pt}\selectfont\renewcommand{\arraystretch}{2}#1}}


\begin{document}

% TODO lo que esté entre estas llaves tendrá tamaño 30pt real
{\fontsize{30pt}{36pt}\selectfont

\begin{center}
    
\textbf{\\CERTAMEN 3:}
\textbf{Econometría ICS-294}
\\
\vspace{0.2cm}
I Semestre 2025\\
Profesores: Jocelyn Tapia Stefanoni, Franco Antonucci Fernández y Ricardo Pruneda Urz\'ua\\
%Ayudante: 
\end{center}

\textbf{Instrucciones:} 

\begin{itemize}

    \item Usted dispone de 90 minutos  para intentar los 100 puntos del certamen.
    \item Cuide la presentación y redacción de sus preguntas.
    \item No puede utilizar celular, computador o tablet, pero sí una calculadora TI.
     
\end{itemize}
\newpage
{\bf RdA3.2  Preguntas cortas (40 puntos)}

{\bf Elija solo 4 preguntas y responda. }
\begin{enumerate}

\item Suponga el siguiente modelo poblacional:
$$PAT = \gamma_0 + \gamma_1 \cdot \text{hombre} + \gamma_2 \cdot econ + \gamma_3 \cdot \text{hombre} \times econ +u$$

Donde: $PAT$ es el patrimonio que acumula una persona; hombre 1 si individuo es hombre, 0 si mujer; $econ = 1$ si estudia economía, 0 si no.
\begin{enumerate}
    \item[(a)] (5 puntos) ¿Cuál es el efecto esperado de haber estudiado economía sobre el patrimonio acumulado?
    \item[(b)] (5 puntos) ¿Difiere este efecto para
hombres y mujeres? Explique.
\end{enumerate}

\newpage

\item 
Se presenta el siguiente modelo econométrico para explicar las importaciones (IMP) en función del Producto Interno Bruto (PIB), la Inversión (INV) y el Consumo (CONS):

\[
\text{IMP} = -19.7251 + 0.0322 \cdot \text{PIB} + 0.4142 \cdot \text{INV} + 0.2427 \cdot \text{CONS}
\]

Considerando la información disponible, responda lo siguiente:

\begin{enumerate}
    \item(6 puntos) ¿Existe evidencia de multicolinealidad entre las variables explicativas? Fundamente su respuesta.

Para su respuesta considere la siguiente información.

    Matriz de correlaciones:

    \[
    \begin{array}{c|ccc}
    & \text{PIB} & \text{INV} & \text{CONS} \\
    \hline
    \text{PIB} & 1 & 0.2154 & 0.9989 \\
    \text{INV} & 0.2154 & 1 & 0.2137 \\
    \text{CONS} & 0.9989 & 0.2137 & 1 \\
    \end{array}
    \]

    Coeficientes de determinación de regresiones auxiliares para cada regresor:

    \[
    R_{PIB}^2 = 0.9979, \quad R_{INV}^2 = 0.0475, \quad R_{CONS}^2 = 0.9979
    \]


    \item (4 puntos) En caso de confirmarse la existencia de multicolinealidad, ¿qué métodos o estrategias podrían emplearse para mitigar este problema en el modelo?

\newpage

    \item (10 puntos)
    ¿Cómo podría utilizar el test de White (simplificado) para detectar heterocedasticidad en el modelo \( y_i = \beta_0 + \beta_1 x_{1i} + \beta_2 x_{2i} + u_i \)? Describa paso a paso el procedimiento e indique cómo se interpreta el resultado.
\end{enumerate}

\item  Sea el proceso estocástico \( x_t \) definido por:
\[
y_t = y_{t-1} + \varepsilon_t,
\]
donde \( \varepsilon_t \) es un proceso de ruido blanco. 

\begin{enumerate}
    \item[(a)] (3 puntos) ¿Puede afirmarse que \( y_t \) es un \textit{random walk}? Justifique conceptualmente.
    
    \item[(b)]  (3 puntos) ¿Es el proceso \( y_t \) estacionario? Justifique su respuesta, considerando tanto la media como la varianza del proceso.
    
    \item[(c)]  (4 puntos) ¿Influye el hecho de que \( \varepsilon_t \) sea gaussiano en la clasificación de \( y_t \) como \textit{random walk}? Explique brevemente.
\end{enumerate}

\item (10 puntos)
En el estudio de Angrist y Krueger (1991), se utiliza el trimestre de nacimiento como variable instrumental para los años de educación en una regresión sobre ingresos. 
 Explique por qué el trimestre de nacimiento podría ser considerado un instrumento válido. 
    Explique qué es el efecto badana.

\newpage
\item Un investigador ha simulado tres series temporales de 200 observaciones. Dos de ellas corresponden a procesos \textbf{AR(1)} con distintos valores del parámetro \( \phi \), y la tercera corresponde a un proceso de \textbf{ruido blanco} (white noise). Para cada caso se presentan dos gráficos:

\begin{itemize}
    \item A la izquierda: la serie temporal simulada \( \{X_t\} \)
    \item A la derecha: su función de autocorrelación (ACF) hasta el rezago 30
\end{itemize}

\newpage
Los valores de \( \phi \) utilizados para los procesos AR(1) fueron:

\[
\phi = -0.95 \quad \text{y} \quad \phi = 0.95
\]

\[
X_t = \phi X_{t-1} + \varepsilon_t
\]

Los gráficos han sido rotulados como \textbf{Caso A}, \textbf{Caso B} y \textbf{Caso C} (sin indicar el valor de \( \phi \) ni el tipo de proceso que representan).

\begin{enumerate}
    \item (4 puntos)  Identifique cuál de los tres casos corresponde al proceso de \textbf{ruido blanco}. Justifique su respuesta con base en el comportamiento de la serie y su ACF.

\newpage
    \item (6 puntos) Para los otros dos casos, asigne a cada uno el valor de \( \phi \) correspondiente. Justifique su elección explicando cómo el signo y la magnitud de \( \phi \) se reflejan en la dinámica de la serie y en la forma de su ACF.

   
\end{enumerate}

\begin{figure}[H]
    \centering
    \includegraphics[width=0.9\linewidth]{certamenes/autocorr.png}
\end{figure}

\end{enumerate}

\newpage
{\bf RdA3.2 Ejercicios (60 puntos)}


{\bf A continuación elija solo 3 preguntas y responda. }

\begin{enumerate}
 \item  Con ayuda de los datos del archivo \texttt{GPA2.RAW} se obtuvo la siguiente ecuación estimada:

\begin{figure}[H]
    \centering
    \includegraphics[width=1\linewidth]{certamenes/reg_c3_ren.1.png}
\end{figure}
\begin{figure}[H]
    \centering
    \includegraphics[width=1\linewidth]{certamenes/reg_c3_ren.2.png}
\end{figure}

\vspace{1em}
\noindent
Otros datos:
\[
n = 4{,}137, \quad R^2 = 0.0858
\]

\newpage

\noindent
Las variables utilizadas son:

\begin{itemize}
    \item \textbf{sat:} Puntuación combinada en el examen de admisión (SAT).
    \item \textbf{hsize:} Tamaño de la cohorte de estudiantes (en centenas) que terminaron el bachillerato en la escuela del estudiante.
    \item \textbf{female:} Variable binaria que vale 1 si el estudiante es mujer y 0 si es hombre.
    \item \textbf{lowincome:} Variable binaria que vale 1 si el estudiante es de bajos ingresos y 0 en caso contrario.
    \item \textbf{female $\times$ lowincome:} Interacción entre género e ingreso bajo.
\end{itemize}

\newpage
\begin{enumerate}
    \item[i)] (5 puntos) ¿Existe alguna evidencia fuerte de que $hsize^2$ deba ser incluida en el modelo? Con base en esta ecuación, ¿cuál es el tamaño óptimo de una escuela?
    
    \item[ii)] (5 puntos) Manteniendo $hsize$ constante, ¿cuál es la diferencia estimada en la puntuación del SAT entre mujeres de ingresos altos y hombres de altos ingresos? ¿Es estadísticamente significativa  esta diferencia?
    \newpage
    \item[iii)] (5 puntos) ¿Cuál es la diferencia estimada en la puntuación del SAT entre hombres de ingresos altos y hombres de ingreso bajo? Pruebe la hipótesis nula de que no hay diferencia entre estas puntuaciones contra la alternativa de que sí hay diferencia.
    
    \item[iv)] (5 puntos) ¿Cuál es la diferencia estimada en la puntuación del SAT entre mujeres de bajos ingresos y mujeres de ingresos altos? ¿Qué se necesita hacer para probar si esta diferencia es estadísticamente significativa?
\end{enumerate}

\newpage
\item Un profesor está analizando qué factores explican el rendimiento académico de los estudiantes en su curso de Econometría. Estima el siguiente modelo de regresión:


\[
\text{nota}_i = \beta_0 + \beta_1 \cdot \text{estudio}_i + \beta_2 \cdot \text{trabaja}_i + u_i
\]

    \newpage


Donde:
\begin{itemize}
    \item $\text{nota}_i$: nota final del estudiante \(i\) en el curso,
    \item $\text{estudio}_i$: horas promedio de estudio a la semana,
    \item $\text{trabaja}_i$: variable dicotómica que vale 1 si el estudiante trabaja durante el semestre, 0 en caso contrario,
    \item $u_i$: término de error, que incluye factores no observados como concentración, salud, o estrés.
\end{itemize}

El profesor estima su modelo y obtiene lo siguiente: 

\textbf{Resultados de R  ** indica 5\% de significancia, *** indica 1\% de significancia}

\begin{figure}[H]
    \centering
    \includegraphics[width=1\linewidth]{certamenes/reg_c3_ren.3.png}
\end{figure}

\begin{enumerate}
    \item (5 puntos) ¿Por qué podría haber heterocedasticidad en este contexto? ¿Por qué representa un problema para la estimación por MCO?

    \item (5 puntos) Plantee la regresión auxiliar del test de White ampliado para este problema.

    \item (10 puntos) Plantee el test de hipótesis de heterocedasticidad. Si el $R^2$ de la regresión auxiliar asciende al valor de 25\%, ¿existe o no heterocedasticidad al 5\% de significancia? Justifique su respuesta.\\

\end{enumerate}

\item A continuación se presenta el gráfico de la serie temporal del tipo de cambio nominal (pesos por dólar) en Chile, correspondiente al período comprendido entre enero de 1983 y diciembre de 2024, con datos de frecuencia mensual. 


\begin{center}
\textbf{Figura 1: Tipos de cambio (pesos por dólar)}
\begin{figure}[H]
    \centering
    \includegraphics[width=0.85\linewidth]{certamenes/BCChChart.png}
\end{figure}
\end{center}


\newpage
Con el objetivo de evaluar si la serie temporal sigue un proceso de caminata aleatoria (random walk), un grupo de investigadores aplica el test de Dickey-Fuller mediante la estimación de dos modelos auxiliares:

\newpage
\begin{figure}[H]
    \centering
    \includegraphics[width=0.95\linewidth]{certamenes/reg_c3_ren.4.png}
\end{figure}

  Testee con un 1\% de significancia si el tipo cambio nominal tiene raíz unitaria. Explique su resultado.

\newpage
\item Un investigador desea estudiar el efecto de las \textbf{horas de entrenamiento semanal} en el \textbf{rendimiento deportivo} de jugadores de básquetbol universitario. Plantea el siguiente modelo:

\[
y_i = \beta_0 + \beta_1 \cdot x_i + \beta_2 \cdot h_i + u_i
\]

Donde:
\begin{itemize}
    \item $y_i$: promedio de puntos anotados por partido durante la temporada.
    \item $x_i$: número de horas de entrenamiento individual por semana. 
    \item $h_i$: altura del jugador en centímetros (variable de control). 
    \item $u_i$: factores no observados que afectan el rendimiento (motivación, talento, disciplina, etc.).
\end{itemize}

Sin embargo, se sospecha que la cantidad de horas de entrenamiento está correlacionada con $u_i$. Por ejemplo, los jugadores más motivados podrían entrenar más y también rendir mejor, incluso si el entrenamiento no fuera la causa del mayor rendimiento.


\begin{enumerate}
    \item (5 puntos) ¿Por qué podría ser inválido estimar el modelo por Mínimos Cuadrados Ordinarios (MCO)? Explique con base en los supuestos del modelo clásico.\\

    \item (5 puntos) Explique por qué $\text{distancia\_cancha}_i$ medida como ``la distancia (en minutos caminando) entre el dormitorio del jugador y la cancha de entrenamiento'' (le llamaremos z) podría ser un \textbf{instrumento válido} para $\text{entreno}_i$. ¿Qué condiciones debe cumplir z para ser un buen instrumento? Exprese matemáticamente.\\

    \newpage
    \item (5 puntos) Suponga que 
    
    $\operatorname{Cov}(x_i, z_i) = -2,94765$ y $\operatorname{Cov}(y_i, z_i) = -4,3215$ 
    
    ¿Qué valor debería tener $\beta_1$? Derive la operación matemática e interprete.\\

    \item (5 puntos) ¿El coeficiente estimado mediante variable instrumental será mayor o menor que el estimado por MCO?\\

\end{enumerate}

\newpage
\item Suponga que usted está interesado en determinar los factores que explican
las diferencias salariales entre distintos individuos. Para realizar
su análisis, usted especifica la siguiente función de regresión:

\begin{equation}
wage = \beta_{1} + \beta_{2}IQ + \beta_{3}Educ + \beta_{4}Exper + \beta_{5}Tenure + \beta_{6}personacolor + u \label{eq:original}
\end{equation}

donde \emph{wage} representa el salario, \emph{IQ} es el coeficiente intelectual,
\emph{Educ} son los años de educación, \emph{Exper} los años de experiencia laboral,
\emph{Tenure} la antigüedad en el trabajo actual, y \emph{personacolor} es una variable dicotómica
que toma el valor 1 cuando la persona se identifica como {persona de color}, y 0 en caso contrario.


Los resultados de la regresión se presentan en la Tabla 1, y en la Figura 1 se muestran los residuos
agrupados según la variable \emph{personacolor}:

\newpage
\begin{center}
\textbf{Tabla 1}
\end{center}

\begin{center}
Model 1: OLS, using observations 1--935\\ Dependent variable: wage\\
\end{center}

\begin{figure}[H]
    \centering
    \includegraphics[width=0.95\linewidth]{certamenes/reg_c3_ren.5.png}
\end{figure}

\begin{center}
\textbf{Figura 1: Residuos agrupados por identificación étnico-racial}
\begin{figure}[H]
    \centering
    \includegraphics[width=0.65\linewidth]{certamenes/racial.png}
\end{figure}
\end{center}

\begin{enumerate}
    \item (4 puntos) De acuerdo a esta información, ¿le parece que existe heterocedasticidad en el modelo? ¿Tiene sentido este resultado? Explique.

 

    \item (7 puntos)  En la Tabla 2 se presentan los resultados de una regresión de los residuos al cuadrado de la regresión de la tabla 1 sobre las variables explicativas:

\newpage
\begin{figure}[H]
    \centering
    \includegraphics[width=0.85\linewidth]{certamenes/reg_c3_ren.6.png}
\end{figure}

\begin{figure}[H]
    \centering
    \includegraphics[width=0.85\linewidth]{certamenes/reg_c3_ren.7.png}
\end{figure}

    \newpage
    ¿Puede usarse esta regresión para testear la presencia de heterocedasticidad? Explique e interprete.

    \item (5 puntos) Observe que la variable \texttt{personacolor} no es estadísticamente significativa según los resultados de la tabla 1. ¿Podemos concluir que \texttt{personacolor} no afecta el ingreso? Fundamente su respuesta considerando la presencia de heterocedasticidad. 

    \item (4 puntos) ¿Cómo se puede corregir el problema anterior para obtener inferencias válidas? ¿Qué tipo de errores estándar se deben utilizar? 

   
\end{enumerate}











\end{enumerate}
\newpage


\begin{center}
    {\Large Formulario  }

 
\end{center}

\begin{itemize}

% \item  
% {\color{red}
% $$\hat{\beta}_0 = \bar{y} - \hat{\beta}_1 \bar{x} $$
% $$\hat{\beta}_1 = \frac{\sum_{i=1}^{n}(x_i - \bar{x})(y_i-\bar{y})}{\sum_{i=1}^{n}(x_i - \bar{x})^2}$$}

\item  {\color{red} \[ \hat{\beta} = (X^{\top}X)^{-1}X^{\top}y \]}

\item 
	{\color{red}	\[ SST = \sum_{i=1}^{n} (y_i - \bar{y})^2 \quad SSE = \sum_{i=1}^{n} (\hat{y}_i - \bar{y})^2  \quad SSR = \sum_{i=1}^{n} \hat{u}_i^2\]}
		

    \item {\color{red} \[ R^2  = 1 - \frac{SSR}{SST} \quad \quad \bar{R}^2 =  1 - \frac{(n-1)SSR}{(n-k-1)SST} \]}

 \item {\color{red} \[ t_{\hat{\beta_j}}=\frac{\hat{\beta}_j - \beta_j}{\text{s.e.}(\hat{\beta}_j)} \sim t_{n-k-1} \]}

 \item {\color{red}\[
    F = \frac{SSR_{r} - SSR_{ur}}{SSR_{ur}} \cdot \frac{n - k - 1}{j} \sim F(j, n - k - 1)
    \]}

\item {\color{red}Test F para regresión auxiliar.}
\begin{center}
 {  \color{red}     \(F=\dfrac{R^{2}/j}{(1-R^{2})/(n-k-1)}
          \sim F(j,\;n-k-1)\).}
\end{center}
% \item {\color{red}Considere que el valor crítico a un nivel de significancia del 5\% es 1.645}
% \item {\color{red}Considere los siguientes valores críticos para t: $\alpha = 0.01 \implies t \approx 2.62$; $\alpha = 0.05 \implies t \approx 1,97$;$\alpha = 0.01 \implies t \approx 2.62$; $\alpha = 0.1 \implies t \approx 1,645$}
%\item {\color{red}Considere que el valor crítico de F al 1\% con grados de libertad (2,178) es aproximadamente 4.72}

\newpage
\item \color{red}{ Multiplicador de Lagrange

$$LM = nR^2 \sim \chi^2_{df}$$

Tabla $\chi^2_{df}$
\begin{center}
    \begin{tabular}{|c|c|c|c|}
    \hline
        df & 0.99 & 0.95 & 0.90 \\ 
        \hline
        1 &  6.635 & 3.841 & 2.706\\
        \hline
        2 &  9.210 & 5.991 & 4.605 \\
        \hline
        3 & 11.345 & 7.815 & 6.251  \\
        \hline
        4 &  13.277 & 9.488 & 7.779 \\
        \hline
        5 &  15.086 & 11.070 & 9.236 \\
        \hline
        6 &  16.812 & 12.592 & 10645 \\
        \hline
    \end{tabular}
    \end{center}}

\newpage
\item Tabla de test D-F
\begin{center}
    \begin{figure}[H]
        \centering
        \includegraphics[width=0.9\linewidth]{sem14/tabla.png}
    \end{figure}
\end{center}

\end{itemize}


% \begin{figure}[H]
%     \centering
%     \includegraphics[width=0.85\linewidth]{Ejercicios//ejercicios_inferencia/cert2_F2_1184.png}
% \end{figure}

% \begin{figure}[H]
%     \centering
%     \includegraphics[width=0.85\linewidth]{certamenes/normal.png}
% \end{figure}


\begin{figure}
    \centering
    \includegraphics[width=0.75\linewidth]{certamenes/normal_1cola.png}
\end{figure}



    \end{document}